\documentclass[a4paper]{article}
\usepackage[polish]{babel}
\usepackage[cp1250]{inputenc}
\usepackage[T1]{fontenc}
\usepackage[dvips]{graphicx}

\pagestyle{headings}
\textwidth      15.5cm
\oddsidemargin    .1cm
\evensidemargin   .1cm

\begin{document}

\thispagestyle{empty}

\begin{minipage}{5cm}
\includegraphics[scale=0.7]{salomon.eps}
\end{minipage}
\begin{minipage}{10cm}
\begin{center}
{\Large Salomon\\
System przetwarzania wiedzy\\}
\end{center}
\end{minipage}

\vspace*{0.5cm}

\hrulefill

\vspace*{1cm}


\begin{flushright}
{\Large \emph{Iteration Description --- E1, ver 1.0} }
\end{flushright}



\begin{flushleft}
{ \Large Historia wersji }\[\]
\begin{tabular}{|p{2.5cm}||p{1.5cm}|p{6cm}|p{3.5cm}|}
  \hline
  Data & Wersja & Opis & Autor \\
  \hline \hline
  29.09.2005 & 1.0 & utworzenie & \emph{Leszek Grzanka}\\
  \hline
  \hline
\end{tabular}
\end{flushleft}

\newpage

\section{Zadania dla iteracji } 
Zaprojektowanie nowej warstwy - \emph{Solution}, reprezentuj�cej zewn�trzne bazy danych. Platforma powinna by� tak�e rozbudowana o mo�liwo�� po��czenia z zewn�trznymi bazami danych.

\section{Produkt iteracji} 
Dokument SAD zawieraj�cy projekt warstwy Solution, oraz interfejsu umo�liwiaj�cego po��czenie z zewn�trznymi bazami danych.

\section{Plan test�w} 
\begin{itemize}
\item Prezentacja projektu warstwy oraz pustych interfejs�w.
\item Konsultacja z drug� grup� pracuj�c� nad Salomonem.
\end{itemize}

\section{Efekt wykonania iteracji}
Efektem by�o opracowanie dokumentu SAD, zawieraj�cego Use-case warstwy \textit{Solution}, model logiczny, implementacyjny, oraz model danych. Wprowadzono tak�e do kodu odpowiednie interfejsy z przeznaczeniem p�niejszego zaimplementowania.

\end{document}

